\section{Related Work}
\label{sec:related}

\textbf{Far-Memory Architecture and Paging Systems.}
Research into remote memory has largely tracked the latency improvements of commodity networks. Early Linux-based systems like Infiniswap \cite{infiniswap} treated remote RDMA memory as an ephemeral block swap device, yielding during page-ins. Because kernel-level page-table operations and interrupt handlers dominated the latency budget, subsequent projects shifted toward aggressive kernel-bypass and busy-waiting. Fastswap \cite{fastswap} streamlined the Linux swap path while spinning for completions, Leap \cite{leap} optimized prefetching heuristics, and Hermit \cite{hermit} used asynchronous worker queues to offload non-critical page-fault processing. DiLOS \cite{dilos} took this approach to its limit by porting the entire network and paging stack into an OSv unikernel, eliminating protection-boundary crossings and delivering microsecond-scale busy-waiting. Adios \cite{adios} builds directly upon the DiLOS unikernel codebase, demonstrating that ultra-cheap unithread context switching makes yielding competitive against DiLOS's spinning model. In contrast to transparent paging, runtime-level disaggregation frameworks such as AIFM \cite{aifm} and Semeru \cite{semeru} require developers to program against remote-memory data structures, avoiding fault-handling penalties at the expense of software transparency.

\textbf{Microsecond Scheduling and Threading Models.}
Handling microsecond-scale I/O has driven significant scheduler innovation. Dataplane operating systems such as ZygOS \cite{zygos} leverage work-stealing across core-pinned queues to mitigate request imbalance. Shinjuku \cite{shinjuku} utilizes hardware virtualization (VMX) timers to preempt long compute requests at 5~$\mu$s intervals, while Concord \cite{concord} relies on compiler-instrumented yield points to avoid hardware timer overhead. Shenango \cite{shenango} and Caladan \cite{caladan} dynamically assign cores to server applications based on packet-arrival congestion, and Persephone \cite{persephone} segregates short and long requests onto dedicated worker pools. However, while these systems treat tasks as compute-bound jobs needing preemption, memory-disaggregated workers stall waiting for network data; preemption cannot accelerate an I/O wait, making thread yielding the appropriate primitive \cite{adios}. Fast user-level context switching has roots in early thread packages such as Capriccio \cite{capriccio} and Scheduler Activations \cite{sa}, later modernized by Arachne \cite{arachne} and Adios's 40-cycle unithreads.

\textbf{Congestion-Aware Load Balancing and Simulation.}
Adios's page-fault-aware dispatcher borrows principles from datacenter load-balancing fabrics like CONGA \cite{conga} and PLB \cite{plb}, using local QP queue depths as congestion signals to steer work away from temporarily backlogged cores. For evaluation, open-loop Poisson arrival generation remains the standard methodology for exposing tail latency bottlenecks without the masking effects of closed-loop client backoff \cite{leverich,tail,tales}. We implement our discrete-event models in SimPy \cite{simpy}.
